% BHU PYQ -- Manually transcribed & error-corrected
% Paper: GLB-603: Hydrogeology, Environmental Geology, Exploration Geology and Computer Application
% Exam: B.Sc. (Hons.) Semester VI Examination, 2022-23

\documentclass[11pt,a4paper]{article}
\usepackage[a4paper,top=2.5cm,bottom=2.5cm,left=2.8cm,right=2.8cm,headheight=14pt]{geometry}
\usepackage{amsmath,amssymb}
\usepackage[T1]{fontenc}
\usepackage[utf8]{inputenc}
\usepackage{lmodern,microtype,fancyhdr,enumitem,hyperref,lastpage,parskip}

\hypersetup{colorlinks=true,urlcolor=black,linkcolor=black,
  pdftitle={GLB-603 Hydrogeology, Environmental Geology, Exploration Geology and Computer Application -- BHU B.Sc. Sem VI 2022-23}}
\pagestyle{fancy}\fancyhf{}
\renewcommand{\headrulewidth}{0.4pt}\renewcommand{\footrulewidth}{0.4pt}
\fancyhead[L]{\small\textit{Banaras Hindu University}}
\fancyhead[R]{\small\textit{GLB-603: Hydrogeology, Env. \& Expl. Geology}}
\fancyfoot[C]{\small Page \thepage\ of \pageref{LastPage}}

\newlist{parts}{enumerate}{2}
\setlist[parts,1]{label=(\alph*),leftmargin=*,itemsep=5pt,topsep=3pt}
\newcommand{\pts}[1]{\hfill\textit{[#1]}}

\begin{document}
\begin{center}
  {\large\bfseries BANARAS HINDU UNIVERSITY}\\[2pt]
  {\normalsize B.Sc. (Hons.) --- Semester VI Examination, 2022-23}\\[6pt]
  \rule{0.8\linewidth}{0.5pt}\\[6pt]
  {\large\bfseries Paper: GLB-603 --- Hydrogeology, Environmental Geology,\\Exploration Geology and Computer Application}\\[4pt]
  {\normalsize Time: 3 Hours \hfill Maximum Marks: 70}\\[2pt]
  {\small Ref. No.: 058442}
\end{center}
\rule{\linewidth}{0.4pt}
\smallskip
\noindent\textit{Instructions: Answer \textbf{five} questions in all, selecting at least \textbf{two} each from Sections A and B, and Section-C which is compulsory. All questions are of equal value.}
\bigskip

\bigskip\noindent{\large\bfseries SECTION --- A}
\rule{\linewidth}{0.4pt}\smallskip

\noindent\textbf{Question 1.} Describe Darcy's law. A confined aquifer is $30\text{ m}$ thick and $7\text{ km}$ wide. Two observation wells are located $1.0\text{ km}$ apart in the direction of flow. The head in well $A$ is $97.0\text{ m}$ and in well $B$ it is $92.0\text{ m}$. If the total daily flow of water through the aquifer is $1050\text{ m}^3\text{/day}$, then calculate the hydraulic conductivity of the aquifer. \pts{14}

\medskip\noindent\textbf{Question 2.} Describe the basic terms related to aquifer properties. Calculate groundwater storage for an aquifer having an areal extent of $10\text{ km}^2$ and saturated thickness $15\text{ m}$. Consider the specific yield of the aquifer is $20\%$. \pts{14}

\medskip\noindent\textbf{Question 3.} Explain the processes and theories used to explain monsoons. Detail their significance in the Indian context as well. \pts{14}

\medskip\noindent\textbf{Question 4.} What are landslides? Explain in detail with their mechanism and classification. What are the global air circulation patterns and how do they affect the climate? \pts{14}

\bigskip\noindent{\large\bfseries SECTION --- B}
\rule{\linewidth}{0.4pt}\smallskip

\noindent\textbf{Question 5.} Write short notes on \textbf{any four} of the following: \pts{$4 \times 3\frac{1}{2} = 14$}
\begin{parts}
  \item Ghyben-Herzberg equation of seawater intrusion
  \item Origin of different types of springs
  \item Hydrochemical properties used to characterize groundwater quality
  \item Thornthwaite classification
  \item Runoff process
\end{parts}

\medskip\noindent\textbf{Question 6.} What are the different exploration methods used for mineral exploration? How are electrical resistivity methods used for exploration? Define the difference between profiling and sounding. \pts{14}

\medskip\noindent\textbf{Question 7.} Explain the fundamental principle of electromagnetic induction with a suitable diagram. What is the self-potential method? Discuss how you acquire SP data and its interpretation. \pts{14}

\bigskip\noindent{\large\bfseries SECTION --- C}
\rule{\linewidth}{0.4pt}\smallskip

\noindent\textbf{Question 8.} Define or explain the following in brief: \pts{$7 \times 2 = 14$}
\begin{parts}
  \item Effluent and influent conditions
  \item Water level and water table
  \item Constant rate of infiltration
  \item Electrical logging
  \item Assay techniques
  \item Global warming
  \item Intensity of earthquakes
\end{parts}

\vfill\rule{\linewidth}{0.4pt}
\begin{center}\small
\href{https://bhu-exam.vercel.app/}{Click here to visit our website}\\[4pt]
\href{https://me-aryan.vercel.app/}{Click here to visit our contributor}\\[4pt]
\href{https://quashan-me.vercel.app/}{Click here to visit our contributor}
\end{center}
\end{document}
